\chapter[27]{Noncompact Gradient Ricci Solitons}
\label{ch:chow-iv-noncompact-gradient-ricci-solitons}
\dcref{ch27}
\source{chow-et-al-ricci-flow-part-iv:page-0022}

Throughout, a gradient Ricci soliton (GRS) is a quadruple
$(\mathcal M^n,g,f,\varepsilon)$ satisfying
\[
\Rc+\nabla^2f+\frac{\varepsilon}{2}g=0.
\]
We use the sign convention of the source: $\varepsilon=-1,0,1$ denotes a
shrinker, steady soliton, or expander, respectively.

\section{Basic properties}

\begin{definition}[GRS structure and normalization]
\label{def:chow-iv-grs-structure}
\source{chow-et-al-ricci-flow-part-iv:page-0023}
\dcref{ch27:eq-27.1--27.6}
\group{Gradient Ricci solitons}

A GRS structure is a connected manifold, metric, potential $f$, and constant
$\varepsilon$ satisfying $\Rc+\nabla^2f+\frac{\varepsilon}{2}g=0$.  It is complete
when $g$ is complete.  The standard identities are
\begin{align*}
R+\Delta f&=-\frac{n\varepsilon}{2},&
\nabla R&=2\Rc(\nabla f),\\
\Delta R+2|\Rc|^2+\varepsilon R&=\langle\nabla R,\nabla f\rangle,&
R+|\nabla f|^2+\varepsilon f&=\mathrm{const}.
\end{align*}
For $\varepsilon=\pm1$ normalize the last constant to zero; for a nonflat
steady soliton normalize it to one.
\end{definition}

\begin{definition}[$f$-Laplacian and weighted curvature]
\label{def:chow-iv-weighted-operators}
\source{chow-et-al-ricci-flow-part-iv:page-0023,chow-et-al-ricci-flow-part-iv:page-0024}
\dcref{ch27:eq-27.7--27.10}
\group{Weighted geometry}

Set $\Delta_f=\Delta-\nabla f\cdot\nabla$, $\Rc_f=\Rc+\nabla^2f$, and
$R_f=R+2\Delta f-|\nabla f|^2$.  The weighted volume is
$\operatorname{Vol}_f(X)=\int_Xe^{-f}\,d\mu$.  For a normalized GRS,
\[
\Delta_f f=\varepsilon\left(f-\frac n2\right)\quad(\varepsilon=\pm1),
\qquad \Delta_f f=-1\quad(\varepsilon=0),
\]
and
\[
\Delta_fR=-2|\Rc|^2-\varepsilon R\le -\frac2nR^2-\varepsilon R.
\]
\end{definition}

\begin{theorem}[Lower bound for scalar curvature of GRS (Theorem 27.2)]
\label{thm:chow-iv-scalar-curvature-lower-bound}
\source{chow-et-al-ricci-flow-part-iv:page-0024,chow-et-al-ricci-flow-part-iv:page-0025,chow-et-al-ricci-flow-part-iv:page-0026,chow-et-al-ricci-flow-part-iv:page-0027}
\uses{def:chow-iv-weighted-operators}
\dcref{ch27:27.2}
\group{Scalar curvature bounds}


If $(\mathcal M^n,g,f,\varepsilon)$ is complete, then $R\ge0$ for a shrinker
or steady soliton, and $R\ge-n/2$ for an expander.
\end{theorem}
\begin{proof}
\sketch
Use a radial cutoff $\Phi_c=\eta(r/c)R$ in the differential inequality for
$\Delta_fR$.  The cutoff satisfies
$\eta''-2(\eta')^2/\eta\ge-\mathrm{const}$; the weighted distance comparison
estimate controls $\Delta_fr$.  The minimum principle gives $R\ge0$ when
$\varepsilon\le0$ and, with a piecewise linear cutoff, $R\ge-n/2$ when
$\varepsilon=1$.
\end{proof}

\begin{remark}[Sharpness (Remark 27.3)]
\label{rem:chow-iv-scalar-lower-sharp}
\source{chow-et-al-ricci-flow-part-iv:page-0024}
\uses{thm:chow-iv-scalar-curvature-lower-bound}
\dcref{ch27:27.3}
\group{Scalar curvature bounds}


The Gaussian soliton realizes $R=0$ in the nonexpanding case, while Einstein
solutions with $\Rc=-g/2$ realize $R=-n/2$ in the expanding case.
\end{remark}

\begin{theorem}[Bounds for the potential gradient (Theorem 27.4)]
\label{thm:chow-iv-gradient-potential-bounds}
\source{chow-et-al-ricci-flow-part-iv:page-0028}
\uses{thm:chow-iv-scalar-curvature-lower-bound}
\dcref{ch27:27.4}
\group{Potential estimates}


For a complete normalized GRS and any base point $\widetilde O$:
\begin{enumerate}
\item a steady soliton satisfies $|\nabla f|\le1$;
\item a shrinker has $f\ge0$ and
\ $|\nabla f|(x)\le\sqrt{f(x)}\le\sqrt{f(\widetilde O)}+d(x,\widetilde O)/2$;
\item an expander has $f\le n/2$ and
\ $|\nabla f|(x)\le\sqrt{n/2-f(x)}\le d(x,\widetilde O)/2+\sqrt{n/2-f(\widetilde O)}$.
\end{enumerate}
\end{theorem}

\begin{remark}[Gaussian sharpness (Remark 27.5)]
\label{rem:chow-iv-gradient-potential-sharp}
\source{chow-et-al-ricci-flow-part-iv:page-0028}
\uses{thm:chow-iv-gradient-potential-bounds}
\dcref{ch27:27.5}
\group{Potential estimates}


For the Gaussian soliton, $|\nabla f|=|\varepsilon|\,|x|/2$, so the preceding
growth rates are qualitatively sharp.
\end{remark}

\begin{lemma}[Criterion for completeness of a vector field (Lemma 27.6)]
\label{lem:chow-iv-vector-field-completeness}
\source{chow-et-al-ricci-flow-part-iv:page-0028}
\dcref{ch27:27.6}
\group{Completeness}

If $F:[0,\infty)\to[0,\infty)$ is locally Lipschitz and every solution of
$u'=F(u)$ is forward complete, then a vector field $X$ on a pointed complete
manifold is complete whenever $|X|(x)\le F(d(x,\widetilde O))$.
\end{lemma}

\begin{corollary}[Completeness of $\nabla f$ (Corollary 27.7)]
\label{cor:chow-iv-gradient-complete}
\source{chow-et-al-ricci-flow-part-iv:page-0029}
\uses{thm:chow-iv-gradient-potential-bounds, lem:chow-iv-vector-field-completeness}
\dcref{ch27:27.7}
\group{Completeness}


Every complete GRS has complete gradient vector field $\nabla f$.
\end{corollary}

\begin{proposition}[Equality case of scalar-curvature lower bound (Proposition 27.8)]
\label{prop:chow-iv-scalar-lower-equality}
\source{chow-et-al-ricci-flow-part-iv:page-0029}
\uses{thm:chow-iv-scalar-curvature-lower-bound, cor:chow-iv-gradient-complete}
\dcref{ch27:27.8}
\group{Equality cases}


For a complete GRS: if $\varepsilon=1$ and $R=-n/2$ somewhere, then
$\Rc=-g/2$; if $\varepsilon=0$ and $R=0$ somewhere, then $\Rc=0$; if
$\varepsilon=-1$ and $R=0$ somewhere, the soliton is Gaussian and isometric to
Euclidean space.
\end{proposition}

\begin{lemma}[A characterization of Euclidean space (Lemma 27.9)]
\label{lem:chow-iv-hessian-characterizes-euclidean}
\source{chow-et-al-ricci-flow-part-iv:page-0030}
\uses{prop:chow-iv-scalar-lower-equality}
\dcref{ch27:27.9}
\group{Equality cases}


If a complete Riemannian manifold admits $f$ with $\nabla^2f=g/2$, then it is
isometric to Euclidean space.  Indeed, after adding a constant,
$f=|\nabla f|^2$, and $r=2\sqrt f$ has $|\nabla r|=1$ and
$\nabla^2(r^2)=2g$; the homothetic level sets are round spheres.
\end{lemma}

\section{Estimates for potential functions}

\begin{corollary}[Bounds on $f$ for GRS (Corollary 27.10)]
\label{cor:chow-iv-potential-bounds}
\source{chow-et-al-ricci-flow-part-iv:page-0031}
\uses{thm:chow-iv-gradient-potential-bounds}
\dcref{ch27:27.10}
\group{Potential estimates}


For a normalized complete GRS and base point $\widetilde O$,
\begin{align*}
\varepsilon=0:&\quad |f(x)-f(\widetilde O)|\le d(x,\widetilde O),\\
\varepsilon=-1:&\quad f(x)\le\tfrac14(d(x,\widetilde O)+2\sqrt{f(\widetilde O)})^2,\\
\varepsilon=1:&\quad -\tfrac14(d(x,\widetilde O)+\sqrt{-4f(\widetilde O)+2n})^2+\tfrac n2
\le f(x)\le\tfrac n2.
\end{align*}
For a shrinker minimum point $O$, $f(O)\le n/2$ and
$f(x)\le\frac14(d(x,O)+\sqrt{2n})^2$.
\end{corollary}

\begin{theorem}[Lower bound for $f$ on shrinkers (Theorem 27.11)]
\label{thm:chow-iv-potential-lower-bound}
\source{chow-et-al-ricci-flow-part-iv:page-0031,chow-et-al-ricci-flow-part-iv:page-0032}
\uses{cor:chow-iv-potential-bounds}
\dcref{ch27:27.11}
\group{Potential estimates}


For a normalized shrinker,
\[
f(x)\ge\frac14\left((d(x,\widetilde O)-2\sqrt{f(\widetilde O)}-4n+\tfrac43)_+\right)^2.
\]
If $O$ minimizes $f$, then
$f(x)\ge\frac14((d(x,O)-35n/8)_+)^2$.
\end{theorem}

\begin{remark}[Comparison with Part III (Remark 27.12)]
\label{rem:chow-iv-potential-part-iii-comparison}
\source{chow-et-al-ricci-flow-part-iv:page-0032}
\dcref{ch27:27.12}
\group{Potential estimates}

The endpoint estimate in the proof is compared with Lemma 19.46 of Part III.
\end{remark}

\begin{corollary}[Uniqueness scale for minima of $f$ (Corollary 27.13)]
\label{cor:chow-iv-potential-minima-distance}
\source{chow-et-al-ricci-flow-part-iv:page-0032}
\uses{thm:chow-iv-potential-lower-bound}
\dcref{ch27:27.13}
\group{Potential estimates}


If $O_1,O_2$ are minimum points of a complete normalized shrinker, then
$d(O_1,O_2)\le35n/8+\sqrt{2n}$.
\end{corollary}

\begin{lemma}[Quadratic upper bound for scalar curvature (Lemma 27.15)]
\label{lem:chow-iv-scalar-quadratic-upper}
\source{chow-et-al-ricci-flow-part-iv:page-0033}
\uses{cor:chow-iv-potential-bounds, thm:chow-iv-scalar-curvature-lower-bound}
\dcref{ch27:27.15}
\group{Scalar curvature bounds}


For a normalized noncompact shrinker,
\[
0\le R(x)=f(x)-|\nabla f|^2(x)\le\frac14(d(x,\widetilde O)+2\sqrt{f(\widetilde O)})^2.
\]
\end{lemma}

\begin{corollary}[Finite topological type under subquadratic $R$ (Corollary 27.16)]
\label{cor:chow-iv-finite-topological-type}
\source{chow-et-al-ricci-flow-part-iv:page-0033}
\uses{thm:chow-iv-potential-lower-bound, lem:chow-iv-scalar-quadratic-upper}
\dcref{ch27:27.16}
\group{Scalar curvature bounds}


If $R(x)\le\alpha(d(x,\widetilde O)+C)^2$ for some $\alpha<1/4$, then the
shrinker has finite topological type; in particular bounded $R$ suffices.
\end{corollary}

\begin{theorem}[Four-dimensional curvature control (Theorem 27.19)]
\label{thm:chow-iv-four-dimensional-curvature}
\source{chow-et-al-ricci-flow-part-iv:page-0034}
\dcref{ch27:27.19}
\group{Scalar curvature bounds}

If a four-dimensional complete noncompact shrinker has bounded scalar curvature,
then $|\Rm|\le C R$ and hence $|\Rm|$ is bounded.
\end{theorem}

\begin{lemma}[Linear-growth scalar curvature on a net (Lemma 27.20)]
\label{lem:chow-iv-linear-net-scalar}
\source{chow-et-al-ricci-flow-part-iv:page-0034,chow-et-al-ricci-flow-part-iv:page-0035}
\dcref{ch27:27.20}
\group{Scalar curvature bounds}

For every $\delta>0$ on a normalized noncompact shrinker, outside a compact set
each $x$ has a point $y\in B_x(\delta)$ satisfying
$R(y)\le C(n,\delta)(d(y,O)+1)$, where $O$ minimizes $f$.
\end{lemma}

\begin{proposition}[Pinched expander curvature decay (Proposition 27.23)]
\label{prop:chow-iv-expander-curvature-decay}
\source{chow-et-al-ricci-flow-part-iv:page-0035}
\dcref{ch27:27.23}
\group{Scalar curvature bounds}

If an expanding GRS has $\Rc\ge\eta Rg$ with $\eta>0$ and $R>0$, then for
every $\widetilde\eta<\eta$,
$R(x)\le C\exp(-\widetilde\eta d(x,\widetilde O)^2)$.
\end{proposition}

\begin{theorem}[K\"ahler expander rigidity (Theorem 27.24)]
\label{thm:chow-iv-kahler-expander-rigidity}
\source{chow-et-al-ricci-flow-part-iv:page-0036}
\uses{prop:chow-iv-expander-curvature-decay}
\dcref{ch27:27.24}
\group{Expander rigidity}


If an expanding K\"ahler GRS satisfies $R=o(d^{-2}(x,\widetilde O))$, then it is
isometric to $\mathbb C^n$.
\end{theorem}

\begin{corollary}[Nonexistence of positively pinched K\"ahler expanders (Corollary 27.25)]
\label{cor:chow-iv-kahler-expander-nonexistence}
\source{chow-et-al-ricci-flow-part-iv:page-0036}
\uses{prop:chow-iv-expander-curvature-decay, thm:chow-iv-kahler-expander-rigidity}
\dcref{ch27:27.25}
\group{Expander rigidity}


There is no expanding K\"ahler GRS of complex dimension at least two with
$\Rc\ge\eta Rg$, $\eta>0$, and $R>0$.
\end{corollary}

\section{Lower bounds for nonflat nonexpanders}

\begin{theorem}[Quadratic lower bound for nonflat shrinkers (Theorem 27.26)]
\label{thm:chow-iv-shrinker-quadratic-lower}
\source{chow-et-al-ricci-flow-part-iv:page-0036,chow-et-al-ricci-flow-part-iv:page-0037}
\uses{def:chow-iv-weighted-operators, thm:chow-iv-potential-lower-bound}
\dcref{ch27:27.26}
\group{Nonflat shrinkers}


Every complete normalized nonflat shrinker satisfies
$R(x)\ge C_0^{-1}d(x,\widetilde O)^{-2}$ outside a compact set.  Thus
$\operatorname{ASCR}(g)>0$ and any asymptotic cone is nonflat.
\end{theorem}

\begin{remark}[Sharpness of quadratic decay (Remark 27.27)]
\label{rem:chow-iv-shrinker-quadratic-sharp}
\source{chow-et-al-ricci-flow-part-iv:page-0037}
\dcref{ch27:27.27}
\group{Nonflat shrinkers}

K\"ahler shrinkers on powers of tautological line bundles have Euclidean volume
growth and quadratic scalar-curvature decay.
\end{remark}

\begin{theorem}[Exponential lower bound for nonflat steadies (Theorem 27.28)]
\label{thm:chow-iv-steady-exponential-lower}
\source{chow-et-al-ricci-flow-part-iv:page-0037,chow-et-al-ricci-flow-part-iv:page-0038}
\dcref{ch27:27.28}
\group{Steady solitons}

If a normalized steady GRS has $f\le0$ and $f\to-\infty$ at infinity, then
\[
R\ge\frac{1}{\sqrt{n/2+2}}e^f\ge c e^{-d(x,\widetilde O)}.
\]
\end{theorem}

\begin{remark}[Cigar soliton (Remark 27.29)]
\label{rem:chow-iv-cigar}
\source{chow-et-al-ricci-flow-part-iv:page-0038}
\dcref{ch27:27.29}
\group{Steady solitons}

For the cigar, $f=-\log(1+x^2+y^2)$ and $R=e^f$.
\end{remark}

\section{Volume growth of shrinking GRS}

\begin{definition}[Asymptotic volume ratio]
\label{def:chow-iv-avr}
\source{chow-et-al-ricci-flow-part-iv:page-0038}
\dcref{ch27:eq-27.69}
\group{Volume growth}

When the limit exists,
\[
\operatorname{AVR}(g)=\lim_{r\to\infty}
\frac{\operatorname{Vol}B_O(r)}{\omega_nr^n}.
\]
For nonnegative Ricci curvature it exists in $[0,1]$ and is base-point
independent.
\end{definition}

\begin{proposition}[Co-area formula (Proposition 27.30)]
\label{prop:chow-iv-coarea}
\source{chow-et-al-ricci-flow-part-iv:page-0038}
\dcref{ch27:27.30}
\group{Volume growth}

For Lipschitz $f$ and integrable (or nonnegative measurable) $h$,
$\int h|\nabla f|\,d\mu=\int dc\int_{\{f=c\}}h\,d\sigma$.
\end{proposition}

\begin{lemma}[Sublevel-set derivatives (Lemma 27.31)]
\label{lem:chow-iv-sublevel-derivatives}
\source{chow-et-al-ricci-flow-part-iv:page-0040}
\uses{prop:chow-iv-coarea}
\dcref{ch27:27.31}
\group{Volume growth}


For $V(c)=\operatorname{Vol}\{f<c\}$ and $\mathcal R(c)=\int_{\{f<c\}}R$,
at regular values $V'(c)=\int_{f=c}|\nabla f|^{-1}$ and
$\mathcal R'(c)=\int_{f=c}R|\nabla f|^{-1}$.  Integration of
$R+\Delta f=n/2$ also gives
\[
\frac n2V(c)-\mathcal R(c)=\int_{f=c}|\nabla f|,
\]
\end{lemma}

\begin{lemma}[Differential identity relating $V$ and $\mathcal R$ (Lemma 27.32)]
\label{lem:chow-iv-sublevel-differential-identity}
\source{chow-et-al-ricci-flow-part-iv:page-0040}
\uses{lem:chow-iv-sublevel-derivatives}
\dcref{ch27:27.32}
\group{Volume growth}


For almost every $c$,
\[
\frac n2V(c)-cV'(c)=\mathcal R(c)-\mathcal R'(c).
\]
\end{lemma}

\begin{theorem}[AVR exists and is dimensionally bounded (Theorem 27.33)]
\label{thm:chow-iv-avr-exists}
\source{chow-et-al-ricci-flow-part-iv:page-0041,chow-et-al-ricci-flow-part-iv:page-0042}
\uses{lem:chow-iv-sublevel-derivatives, lem:chow-iv-sublevel-differential-identity, def:chow-iv-avr}
\dcref{ch27:27.33}
\group{Volume growth}


For every complete normalized noncompact shrinker, $\operatorname{AVR}(g)$ exists
and is bounded above by a constant depending only on $n$.  With
$P(c)=V(c)c^{-n/2}-\mathcal R(c)c^{-n/2-1}$,
$P'(c)\le0$ for $c\ge(n+2)/2$ and
$2^n\omega_n\operatorname{AVR}(g)=\lim_{c\to\infty}P(c)$.
\end{theorem}

\begin{proposition}[Characterization of positive AVR (Proposition 27.35)]
\label{prop:chow-iv-positive-avr}
\source{chow-et-al-ricci-flow-part-iv:page-0042}
\uses{thm:chow-iv-avr-exists}
\dcref{ch27:27.35}
\group{Volume growth}


For a complete normalized noncompact shrinker,
\[
\operatorname{AVR}(g)>0\iff
\int_{n+2}^{\infty}\frac{\mathcal R(c)}{cV(c)}\,dc<\infty
\iff
\int_1^\infty\frac{\int_{B_O(r)}R}{\operatorname{Vol}B_O(r)}\frac{dr}{r}<\infty.
\]
\end{proposition}

\begin{corollary}[Volume growth with a scalar lower bound (Corollary 27.36)]
\label{cor:chow-iv-volume-lower-scalar}
\source{chow-et-al-ricci-flow-part-iv:page-0043}
\uses{prop:chow-iv-positive-avr}
\dcref{ch27:27.36}
\group{Volume growth}


If $R\ge\delta\ge0$, then
$\operatorname{Vol}B_O(r)\le C(\mathcal G,O)(1+r)^{n-2\delta}$.
\end{corollary}

\begin{remark}[Cylinder shrinkers (Remark 27.37)]
\label{rem:chow-iv-cylinder-shrinkers}
\source{chow-et-al-ricci-flow-part-iv:page-0043}
\dcref{ch27:27.37}
\group{Volume growth}

The products $(\mathcal N^k,h)\times\mathbb R^{n-k}$ with $\Rc_h=h/2$ have
$R=k/2$ and volume growth of order $(1+r)^{n-k}$.
\end{remark}

\begin{theorem}[Bakry--Emery volume comparison (Theorem 27.41)]
\label{thm:chow-iv-bakry-emery-comparison}
\source{chow-et-al-ricci-flow-part-iv:page-0043,chow-et-al-ricci-flow-part-iv:page-0045}
\dcref{ch27:27.41}
\group{Volume growth}

If $\Rc_f\ge0$ and $|\nabla f|\le A$ on $B_O(r)$, then
\[
\operatorname{Vol}_fB_O(r)\le\omega_n e^{-f(O)+Ar}r^n,
\qquad
\operatorname{Vol}B_O(r)\le\omega_ne^{-f(O)+Ar+C}r^n
\]
when $f\le C$ on the ball.
\end{theorem}

\begin{theorem}[At most Euclidean volume growth (Theorem 27.42)]
\label{thm:chow-iv-euclidean-volume-growth}
\source{chow-et-al-ricci-flow-part-iv:page-0045,chow-et-al-ricci-flow-part-iv:page-0046}
\uses{thm:chow-iv-bakry-emery-comparison}
\dcref{ch27:27.42}
\group{Volume growth}


Every complete normalized noncompact shrinker satisfies
\[
\operatorname{Vol}B_O(r)\le\omega_ne^{f(O)}r^n.
\]
If $O$ minimizes $f$, the factor is at most $e^{n/2}$.
\end{theorem}

\begin{remark}[Sharper volume bounds at a scalar minimum (Remark 27.44)]
\label{rem:chow-iv-volume-at-scalar-minimum}
\source{chow-et-al-ricci-flow-part-iv:page-0046}
\dcref{ch27:27.44}
\group{Volume growth}

If $O$ minimizes both $R$ and $f$, then $\operatorname{Vol}B_O(r)\le\omega_nr^n$.
\end{remark}

\section{Logarithmic Sobolev inequality}

\begin{definition}[Perelman entropy and shrinker entropy]
\label{def:chow-iv-shrinker-entropy}
\source{chow-et-al-ricci-flow-part-iv:page-0047}
\dcref{ch27:eq-27.115--27.118}
\group{Logarithmic Sobolev inequality}

For $\int e^{-\varphi}=(4\pi)^{n/2}$, set
\[
\mathcal W(g,\varphi)=\int(R+|\nabla\varphi|^2+\varphi-n)(4\pi)^{-n/2}e^{-\varphi}d\mu,
\quad \mu(g)=\inf_\varphi\mathcal W(g,\varphi).
\]
Normalize a shrinker by $\int e^{-f}=(4\pi)^{n/2}$ and
$R+|\nabla f|^2-f=C_1(\mathcal G)$; then $\mu(\mathcal G)=-C_1(\mathcal G)$.
\end{definition}

\begin{theorem}[Sharp logarithmic Sobolev inequality (Theorem 27.46)]
\label{thm:chow-iv-sharp-log-sobolev}
\source{chow-et-al-ricci-flow-part-iv:page-0047,chow-et-al-ricci-flow-part-iv:page-0048,chow-et-al-ricci-flow-part-iv:page-0049}
\uses{def:chow-iv-shrinker-entropy}
\dcref{ch27:27.46}
\group{Logarithmic Sobolev inequality}


For every complete noncompact shrinker,
\[
\inf_\varphi\mathcal W(g,\varphi)=\mu(\mathcal G)=-C_1(\mathcal G).
\]
Equivalently, if $\int w^2=(4\pi)^{n/2}$,
\[
\int(Rw^2+4|\nabla w|^2-w^2\log w^2)\,d\mu
\ge(n-C_1(\mathcal G))\int w^2d\mu.
\]
\end{theorem}

\begin{remark}[Closed comparison (Remark 27.47)]
\label{rem:chow-iv-log-sobolev-closed}
\source{chow-et-al-ricci-flow-part-iv:page-0048}
\dcref{ch27:27.47}
\group{Logarithmic Sobolev inequality}

The corresponding closed-manifold inequality bounds
$\int(b|\nabla\varphi|^2+\varphi-n)e^{-\varphi}$ below by a metric-dependent constant.
\end{remark}

\begin{lemma}[Fisher and entropy dissipation]
\label{lem:chow-iv-fisher-entropy-dissipation}
\source{chow-et-al-ricci-flow-part-iv:page-0048,chow-et-al-ricci-flow-part-iv:page-0049}
\uses{def:chow-iv-weighted-operators}
\dcref{ch27:eq-27.119--27.127}
\group{Logarithmic Sobolev inequality}


For $\square_f\xi=|\nabla\xi|^2$ and normalized
$\int e^{\xi-f}=(4\pi)^{n/2}$, define
$I(\xi)=\int|\nabla\xi|^2e^{\xi-f}$ and
$H(\xi)=\int\xi e^{\xi-f}$.  Then
\[
I'(t)=-\int e^{\xi}(2|\nabla^2\xi|^2+|\nabla\xi|^2)e^{-f}\le-I(t),
\qquad H'(t)=-I(t),
\]
and hence $H(\xi)\le I(\xi)$, yielding the sharp inequality above.
\end{lemma}

\begin{remark}[Bakry--Emery method (Remark 27.48)]
\label{rem:chow-iv-bakry-emery-method}
\source{chow-et-al-ricci-flow-part-iv:page-0049}
\dcref{ch27:27.48}
\group{Logarithmic Sobolev inequality}

The dissipation computation parallels Bakry--Emery's proof of the classical
logarithmic Sobolev inequality.
\end{remark}

\begin{corollary}[Noncollapsing (Corollary 27.49)]
\label{cor:chow-iv-shrinker-noncollapsing}
\source{chow-et-al-ricci-flow-part-iv:page-0050}
\uses{thm:chow-iv-sharp-log-sobolev}
\dcref{ch27:27.49}
\group{Shrinker geometry at infinity}


A complete noncompact shrinker is $\kappa$-noncollapsed at scales where
$R\le r^{-2}$, with $\kappa$ depending only on its entropy constant.
\end{corollary}

\begin{theorem}[At least linear volume growth (Theorem 27.50)]
\label{thm:chow-iv-shrinker-linear-growth}
\source{chow-et-al-ricci-flow-part-iv:page-0050}
\uses{thm:chow-iv-sharp-log-sobolev}
\dcref{ch27:27.50}
\group{Shrinker geometry at infinity}


For every base point $p$ there is $c>0$ such that
$\operatorname{Area}S(p,r)\ge2c$ for $r\ge1/2$ and
$\operatorname{Vol}B_p(r)\ge cr$ for $r\ge1$.
\end{theorem}

\begin{theorem}[Uniqueness of asymptotic cones (Theorem 27.51)]
\label{thm:chow-iv-asymptotic-cone-uniqueness}
\source{chow-et-al-ricci-flow-part-iv:page-0050}
\dcref{ch27:27.51}
\group{Shrinker geometry at infinity}

Two complete noncompact shrinkers with two ends asymptotic to the same smooth
Euclidean cone have isometric universal covers.
\end{theorem}

\begin{theorem}[Ricci decay implies a Euclidean cone (Theorem 27.52)]
\label{thm:chow-iv-asymptotic-cone}
\source{chow-et-al-ricci-flow-part-iv:page-0050}
\dcref{ch27:27.52}
\group{Shrinker geometry at infinity}

If $|\Rc|(x)\to0$ at infinity on a complete noncompact shrinker, then the
metric is asymptotic to a Euclidean cone over a smooth closed manifold.
\end{theorem}

\section{Gradient shrinkers with nonnegative Ricci curvature}

\begin{theorem}[Positive scalar lower bound for nonflat $\Rc\ge0$ shrinkers (Theorem 27.53)]
\label{thm:chow-iv-positive-scalar-shrinker}
\source{chow-et-al-ricci-flow-part-iv:page-0050,chow-et-al-ricci-flow-part-iv:page-0051,chow-et-al-ricci-flow-part-iv:page-0052}
\uses{cor:chow-iv-gradient-complete, prop:chow-iv-scalar-lower-equality}
\dcref{ch27:27.53}
\group{Nonnegative Ricci curvature}


If a complete noncompact shrinker has $\Rc\ge0$ and is not Ricci flat, then
there is $\delta>0$ with $R\ge\delta$ everywhere.  Along integral curves of
$\nabla f$, $R$ is nondecreasing; the proof uses the estimate
$\langle\nabla f,\gamma'(r)\rangle\ge r/4-C_1-|\nabla f|(\widetilde O)$
and a backward-flow contradiction.
\end{theorem}

\begin{theorem}[Nonnegative Ricci curvature forces zero AVR (Theorem 27.54)]
\label{thm:chow-iv-nonnegative-ricci-avr-zero}
\source{chow-et-al-ricci-flow-part-iv:page-0053}
\uses{thm:chow-iv-positive-scalar-shrinker, cor:chow-iv-volume-lower-scalar}
\dcref{ch27:27.54}
\group{Nonnegative Ricci curvature}


A complete noncompact non-Ricci-flat shrinker with $\Rc\ge0$ has
$\operatorname{AVR}(g)=0$.
\end{theorem}

\begin{remark}[Ricci-flat examples with positive AVR (Remark 27.55)]
\label{rem:chow-iv-ricci-flat-positive-avr}
\source{chow-et-al-ricci-flow-part-iv:page-0053}
\dcref{ch27:27.55}
\group{Nonnegative Ricci curvature}

There are complete noncompact four-dimensional Ricci-flat (hyper-K\"ahler and
asymptotically locally Euclidean) manifolds with positive AVR.
\end{remark}

\begin{conjecture}[Optimistic curvature conjecture (Conjecture 27.18)]
\label{conj:chow-iv-bounded-curvature-conjecture}
\dcref{ch27:27.18}
\group{Open problems}
\source{chow-et-al-ricci-flow-part-iv:page-0034}
Every complete noncompact shrinker has scalar, Ricci, and sectional curvature
bounded above.
\end{conjecture}

\begin{conjecture}[Elliptic Harnack estimate (Conjecture 27.21)]
\label{conj:chow-iv-elliptic-harnack}
\dcref{ch27:27.21}
\group{Open problems}
\source{chow-et-al-ricci-flow-part-iv:page-0035}
There is a universal constant such that $d(x,y)\le1$ implies
$R(x)\le\mathrm{const}\,R(y)$ on every complete noncompact shrinker.
\end{conjecture}

\begin{claim}[Harnack implies finite topological type (Exercise 27.22)]
\label{ex:chow-iv-harnack-topology}
\dcref{ch27:27.22}
\group{Open problems}
\source{chow-et-al-ricci-flow-part-iv:page-0035}
\uses{conj:chow-iv-elliptic-harnack,cor:chow-iv-finite-topological-type}
Show that the elliptic Harnack estimate implies the hypothesis needed for finite
topological type.
\end{claim}

\begin{claim}[Monotone radial density (Exercise 27.43)]
\label{ex:chow-iv-radial-density-monotone}
\dcref{ch27:27.43}
\group{Volume growth}
\source{chow-et-al-ricci-flow-part-iv:page-0046}
\uses{thm:chow-iv-euclidean-volume-growth}
Show that $\left(J(r)/(e^{f(O)}r^{n-1})\right)^r$ is nonincreasing and tends to
one as $r\downarrow0$.
\end{claim}

\begin{claim}[Second monotone density (Exercise 27.45)]
\label{ex:chow-iv-second-radial-density}
\dcref{ch27:27.45}
\group{Volume growth}
\source{chow-et-al-ricci-flow-part-iv:page-0046}
The expression
\[
\frac{e^{r^2/12-f(r)+(2/r)\int_0^rf}}{r^{n-1}}J(r)
\]
is nonincreasing and tends to $e^{f(O)}$ at the origin.
\end{claim}

\begin{claim}[Conjugation of the weighted Schrodinger operator (Exercise 27.1)]
\label{ex:chow-iv-weighted-conjugation}
\dcref{ch27:27.1}
\group{Weighted geometry}
\source{chow-et-al-ricci-flow-part-iv:page-0024}
\uses{def:chow-iv-weighted-operators}
For every smooth $\varphi$,
\[
\left(\Delta_f-\frac14R_f\right)\varphi
=e^{f/2}\left(\Delta-\frac14R\right)(e^{-f/2}\varphi).
\]
\end{claim}

\begin{claim}[AVR upper bound one (Problem 27.34)]
\label{prob:chow-iv-avr-one}
\dcref{ch27:27.34}
\group{Open problems}
\source{chow-et-al-ricci-flow-part-iv:page-0042}
Must every shrinker satisfy $\operatorname{AVR}(g)\le1$?
\end{claim}

\begin{claim}[Critical quadratic threshold (Problem 27.17)]
\label{prob:chow-iv-critical-growth}
\dcref{ch27:27.17}
\group{Open problems}
\source{chow-et-al-ricci-flow-part-iv:page-0034}
Can the hypothesis $\alpha<1/4$ in Corollary~\ref{cor:chow-iv-finite-topological-type}
be removed?
\end{claim}

\begin{claim}[Potential functions of expanders (Problem 27.14)]
\label{prob:chow-iv-expander-potential}
\dcref{ch27:27.14}
\group{Open problems}
\source{chow-et-al-ricci-flow-part-iv:page-0033}
What is the best qualitative estimate for $f$ on expanders?  With $\Rc>0$ one
has the quadratic upper bound $f\le-d(x,\widetilde O)^2/4+C$.
\end{claim}

\begin{claim}[Splitting problem for nonnegative-Ricci shrinkers (Problem 27.56)]
\label{prob:chow-iv-shrinker-splitting}
\dcref{ch27:27.56}
\group{Open problems}
\source{chow-et-al-ricci-flow-part-iv:page-0053}
Does every noncompact shrinker with $\Rc\ge0$ split as a product of a compact
shrinker and Euclidean space?
\end{claim}

\begin{claim}[Compact factors (Question 27.38)]
\label{ques:chow-iv-compact-factor}
\dcref{ch27:27.38}
\group{Open problems}
\source{chow-et-al-ricci-flow-part-iv:page-0043}
Does a positive scalar lower bound force a compact factor of the corresponding
dimension on a simply connected shrinker?
\end{claim}

\begin{claim}[Euclidean growth without a line factor (Question 27.40)]
\label{ques:chow-iv-euclidean-growth-no-line}
\dcref{ch27:27.40}
\group{Open problems}
\source{chow-et-al-ricci-flow-part-iv:page-0043}
Does a simply connected shrinker without an $\mathbb R$ factor have Euclidean
volume growth?
\end{claim}

\begin{claim}[Positive AVR and finite ASCR (Question 27.39)]
\label{ques:chow-iv-positive-avr-finite-ascr}
\dcref{ch27:27.39}
\group{Open problems}
\source{chow-et-al-ricci-flow-part-iv:page-0043}
Does positive AVR imply finite ASCR for a complete noncompact shrinker?
\end{claim}
